\documentclass[UTF8,a4paper,12pt]{ctexart}

\usepackage[a4paper,margin=2.5cm]{geometry}
\usepackage{array}
\usepackage{amsmath}
\usepackage{booktabs}
\usepackage{longtable}
\usepackage{tabularx}
\usepackage{enumitem}
\usepackage{xcolor}
\usepackage{tikz}
\usepackage{eso-pic}
\usepackage{hyperref}
\usepackage{fancyhdr}
\usepackage{lastpage}

\usetikzlibrary{calc}

\hypersetup{
  colorlinks=true,
  linkcolor=black,
  urlcolor=blue
}

\pagestyle{fancy}
\fancyhf{}
\cfoot{\thepage/\pageref{LastPage}}
\renewcommand{\headrulewidth}{0pt}

\setlength{\parindent}{2em}
\setlength{\parskip}{0.35em}
\setlength{\headheight}{15pt}
\setlist[enumerate]{itemsep=0.2em,topsep=0.4em}
\renewcommand{\arraystretch}{1.25}

\newcommand{\schoolname}{北\quad 京\quad 邮\quad 电\quad 大\quad 学}
\newcommand{\reporttitle}{实\quad 验\quad 报\quad 告}
\newcommand{\coursename}{计算机组成原理实验}
\newcommand{\experimentname}{中断原理与中断返回测试}
\newcommand{\collegename}{计算机学院}
\newcommand{\classname}{2024211310}
\newcommand{\studentname}{王皓轩}
\newcommand{\studentid}{2024211213}
\newcommand{\teachername}{石劭轩}
\newcommand{\reportdate}{2026 年 5 月 27 日}

\newcommand{\covervalue}[4]{%
  \noindent\makebox[#1][l]{\bfseries #2}%
  \underline{\makebox[#3][l]{\bfseries\strut #4}}%
}
\newcommand{\coverfield}[1]{\underline{\makebox[4.15cm][l]{\bfseries\strut #1}}}

\AddToShipoutPictureBG{%
  \begin{tikzpicture}[remember picture,overlay]
    \draw[line width=0.8pt]
      ($(current page.north west)+(1.25cm,-1.25cm)$)
      rectangle
      ($(current page.south east)+(-1.25cm,1.25cm)$);
  \end{tikzpicture}%
}

\begin{document}
\hypersetup{pageanchor=false}
\begin{titlepage}
\thispagestyle{empty}
\centering
\vspace*{1.9cm}
{\zihao{2}\bfseries \schoolname\par}
\vspace{1.35cm}
{\zihao{-0}\bfseries \reporttitle\par}

\vspace{3.7cm}
\begin{minipage}{12.8cm}
{\zihao{4}
\covervalue{2.5cm}{课程名称}{9.8cm}{\coursename}\par
\vspace{1.25cm}
\covervalue{2.5cm}{实验名称}{9.8cm}{\experimentname}\par
\vspace{1.25cm}
\begin{tabular}{@{}>{\bfseries}p{1.4cm}@{}p{4.15cm}@{\hspace{0.85cm}}>{\bfseries}p{1.4cm}@{}p{4.15cm}@{}}
学院 & \coverfield{\collegename} & 班级 & \coverfield{\classname} \\
\end{tabular}\par
\vspace{1.1cm}
\begin{tabular}{@{}>{\bfseries}p{1.4cm}@{}p{4.15cm}@{\hspace{0.85cm}}>{\bfseries}p{1.4cm}@{}p{4.15cm}@{}}
姓名 & \coverfield{\studentname} & 学号 & \coverfield{\studentid} \\
\end{tabular}\par
\vspace{1.1cm}
\begin{tabular}{@{}>{\bfseries}p{1.4cm}@{}p{4.15cm}@{\hspace{0.85cm}}>{\bfseries}p{1.4cm}@{}p{4.15cm}@{}}
教师 & \coverfield{\teachername} & 成绩 & \coverfield{} \\
\end{tabular}\par
}
\end{minipage}

\vfill
{\zihao{3}\bfseries \reportdate\par}
\vspace*{1.7cm}
\end{titlepage}
\hypersetup{pageanchor=true}

\section{实验目的}

在此填写实验目的。

\section{实验原理}

在此填写实验原理。

\section{实验步骤}

在此填写实验步骤。

\section{实验数据与分析}

在此填写实验数据与分析。

\section{实验结论}

在此填写实验结论。

\end{document}
